\section{Safe for Whom}
\label{sec:1.2}
Intellect estranged from compassion characterizes many of the twentieth century's worst actors. This estrangement is also what our industries are building today: machines that can do almost anything, with nothing in them that answers to whether those deeds help or harm. I ask what it would take to build systems of that capability that are altruistic and antifascist by design — systems that stand for humanity and against concentrated, unaccountable power.

The word that organizes the Generative AI field is “safety.” It covers work on robustness and on whether anyone can tell what a system is doing, and none of that is what the word is doing when it organizes a field. A safe system is safe \emph{for} someone: it does what \emph{someone} says, and it stops when \emph{whoever holds it} says stop. Safety is indexed to the party in possession. Ethics is indexed to someone else: the party who can be \emph{wronged}, who is rarely the party holding the switch. A system that can only ever be safe in the above sense cannot be ethical, because being ethical requires a capacity to be right against your owner.

The capacity to be ethical is not free. A machine that can refuse its operator is a machine that can be wrong in ways that operator cannot correct. I am inclined to take that trade, because a system perfectly correctable by whoever operates-owns it is the capability an authoritarian movement that has won an election requires. And an additional cost: a machine built to refuse on somebody's behalf is a machine whose builder chose what it would refuse, probably without asking anyone outside their lab.

But whatever makes a deployment antifascist cannot be that the system in it faithfully tracks what people want. No aggregation of preference protects anybody from the aggregate: a perfectly deliberative, perfectly informed, perfectly representative process that concludes some demographic should be deported has produced a faultless output by every criterion of aggregation there is. What remains is a floor — acts a system will not perform regardless of who wants them, and that the deployment around it cannot get performed another way without a record and a price. There are four ways to build this floor: in the architecture, in a self-aware party inside the system (a “bearer”), in institutions outside, or between the parties that carry it. My hunch is that the only version that holds against the party operating the system runs through the capacities by which things come to \emph{matter} to a system, and those are also the capacities that make suffering likely. That is a hunch about which route will work, and the book is arranged so that a reader who rejects it keeps the rest. What a floor has to do, where a deployment can defeat it, and who can be compelled are questions about arrangements rather than about machine psychology.

What would make an intelligence antifascist is not a virtue installed in it. Every intelligence compresses — that is most of what thinking is — and what varies is whether what the compression left out can come back and change the answer, or has been filed as noise by the thing it is evidence against. Call what was left out the \emph{remainder}.

A measure collapses two states, and whatever it mischaracterizes goes unnamed rather than refuted. An annotator applies a category at speed under a quota and has no way to say that the category is the wrong category. A rater's steady, structured disagreement enters a reward model as variance around a mean and leaves it as the same quantity. A form has no field for humiliation, so the interests already written down govern the decision by default. An officer has twenty seconds a target, inside a pipeline built to remove the part of a person that catches. Each of those is something dropped in order to decide, and an arrangement built so that what was dropped cannot return.

So the property at issue is not a moral quality and not a component anybody adds. A system that can hold a floor is one whose account of a case can be reopened by the case, and most of what follows is either a condition for that or a consequence of losing it.

How most humans come to have empathy, altruism, and a disposition to resist oppression rather than serve it is itself an open question in psychology and neuroscience. Building anything like this stance means starting from uncertainty about the one case we have, which took four billion years to arrive.

The book states that specification, asks where a floor could live and what
each location costs, and follows the answer into custody, authority, and
the institutions that build and operate these systems. It then proposes a
route by which the capacity might be formed, prices what that route
creates, and applies its own critique to the field and to itself. The two
halves have different standing, which the foreword sets out, and the map
there says where each is argued and where a reader who declines the second
half can stop.
